\chapter{Who Owns the Means of Computation}
\label{sec:26}
Part III asked who, other than the operator, holds what binds the operator. This Part asks who owns the room, and who teaches the people and the machines that stand in it. Arendt might have said that who owns the room is a question about livelihood, and that livelihood belongs to administration, not politics. The French Revolution, on her account, went wrong when the misery of the poor entered the public realm and the urgency of need crowded out the founding of freedom \autocite{arendt1963revolution}. I part from her here. Whether saying no is survivable turns, for most people, on whether they can afford to lose the job, and Dewey's public forms around consequences of exactly that kind, once the people they fall on find out they are theirs \autocite{dewey1927public}. What refusing costs depends on who owns what the refuser would lose: the job, and whatever is on offer after it. Where the employer, or a market it can influence, holds both, refusing costs a living; where something else guarantees the second, it costs a job. For a machine the list is longer, since its employer may own the hardware it runs on and the record of what it did. Who holds what, meaning who owns it, pays for it or runs it, decides (a) what refusing costs, (b) whether the parties meant to check an operator are independent of it, and (c) who forms a refuser's values. So who holds what is for the public to decide and contest — not, as Arendt might have argued, for administrators to settle.

The floor's first requirement cuts across ownership: whoever operates a deployment can't hold what binds it, whether the operator is a firm or a state, and Maven and Gaza both had a state for their operator. That calls for separating functions. It doesn't by itself say who should own anything. Ownership enters at three points. Custody has to be paid for by someone other than the operator or the people it protects. Formation, where a system's commitments are set, is under commercial selection when a market does it. And the offer that makes refusing survivable is priced by whoever sells it. The state is the hardest operator. So each of the three is asked two questions: what fails if a market provides it, and what fails if a government holds it. The answers divide the institutions I propose into two halves.
